\documentclass[12pt]{article}
\usepackage{amsmath, amsthm, amssymb}
\usepackage[utf8]{inputenc}
\usepackage[margin=1in]{geometry}

\newtheorem{theorem}{Theorem}

\begin{document}

\section{Infinite Unions of Points and Intervals}
\begin{theorem}
Let $\mathcal{M} = (M, <, \dots)$ be an o-minimal structure.
Let $S \subseteq M$ be a set constructed from the union of a finite number of points and open intervals.
If $S$ is infinite, then $S$ must contain an open interval.
\end{theorem}
\begin{proof}
Let the set $S$ be denoted by $S = P \cup I_1 \cup \dots \cup I_k$, where $P$ is a finite set of points and each $I_j$ is an open interval in $M$.
Assume, for the sake of contradiction, that $S$ does not contain an open interval.
This implies that the number of intervals $k = 0$. Consequently, $S = P$.
By definition, $P$ is a finite set of points, meaning its cardinality $|P|$ is finite. Thus, $S$ would be finite.
This contradicts the hypothesis that $S$ is an infinite set.
Therefore, our assumption that $S$ contains no open intervals must be false.
Hence, $k \ge 1$, and $S$ contains at least one open interval.
Furthermore, because the underlying order of an o-minimal structure is dense, any such open interval $(a, b)$ with $a < b$ is strictly infinite, which consistently accounts for the infinite cardinality of $S$. 
Note that an interval extending to positive infinity, yielding $+\infty$ as a supremum, is also trivially an infinite open interval.
\end{proof}

\end{document}
